\SourceLeaf{286}%
\noindent wo $F$ die Fläche der ganzen Figur, $S$ die der ganzen Fläche bedeutet,
\[
= (2\chi - \sigma) \, . \, 360^0 - \frac{F}{S} \,.\,720^0.
\]
Diese Vorschrift begreift die oben gegebene unter sich, da der unendliche Raum in der Ebene den Character 0 hat, wodurch $\frac{F}{S} = 0$ wird; oben wurden die Zeichen da gezählt, wo die Intersectionspunkte zum erstenmale getroffen wurden, und wo sie den hier gebrauchten entgegengesetzt sind.

\begin{center}
[7.]
\end{center}

\begin{center}\textbf{EINE ANDERE ART, DIESE GEGENSTÄNDE ZU BEHANDELN.}\end{center}

Es sei $S$ die ganze Curve, und $s$ ein unbestimmtes Stück. Im Anfangspunkte ziehe man eine Tangente und bezeichne durch $(0, s)$ den Winkel, welchen die gerade Linie vom Anfangspunkte zum Endpunkte des Stücks $s$ macht. Offenbar wird sich $(0, s)$ von $s = 0$ an so lange nach dem Gesetz der Stetigkeit ändern, als die Curve nicht wieder in den Anfangspunkt zurückgekommen ist; wir nehmen an, dass diess nicht eher geschehe, als bis $s = S$ geworden ist, und der letzte Werth, ehe $s$ diesen Werth erreicht hat, wird sein
\[
(0, s) = \pm 180^0,
\]
in so fern man die Winkel nach der innern oder äussern Seite der Linie wachsen lässt.

Indem man $s$ als constant ansieht, ziehe man von jedem Punkt der Linie, die als Endpunkt des kleinern Stücks $s'$ angesehen wird, eine gerade Linie zum Endpunkt des Stücks $s$ und nenne deren Neigung gegen die erste Tangente $(s', s)$. So lange der Endpunkt von $s'$ mit dem von $s$ nicht zusammenfällt, wird sich also, bei stetig geändertem $s'$, $(s', s)$ nach dem Gesetz der Stetigkeit ändern, falls das Stück $s$ sich nicht selbst schneidet, bis $s' = s$. Für diesen Grenzwerth selbst ist $(s', s)$, genau genommen, unbestimmt; so lange aber $s' < s$, wie nahe sich auch die Werthe liegen, ist er bestimmt; und er nähert sich, je näher $s'$ dem $s$ kommt, immer mehr der Neigung der Tangente gegen die Tangente im Anfangspunkte. Es sei nun $s = S - \omega$, und $s'$ verändere sich von $0$ bis $S - \omega - \omega'$; es ist klar, dass die Veränderung von $s'$ dadurch $180^0$ desto näher kommen wird, je kleiner $\omega, \omega'$ sind, so dass $(S - \omega - \omega', S - \omega)$ dem
